% =====================================================================
% 第 4 章  研究设计
% ---------------------------------------------------------------------
% 编译说明同前：ctexart 下用 \section，ctexbook/report 下改为 \chapter。
% 依赖宏包：booktabs、threeparttable、amsmath
% =====================================================================

\section{研究设计}

本章说明识别策略。核心问题是：报告级情绪每期只有一个取值，如何在行业面板中识别它的
影响；以及在这一约束下，本文选择什么样的估计方式、做了哪些检验来支撑结论。

\subsection{识别上的一个约束}

设 $i$ 表示行业、$t$ 表示季度，考虑常见的双向固定效应设定
\begin{equation}
y_{i,t}=\alpha+\beta\,\text{tone}_t+\gamma'X_{i,t}+\mu_i+\lambda_t+\varepsilon_{i,t},
\label{eq:twoway}
\end{equation}
其中 $\mu_i$ 为行业固定效应，$\lambda_t$ 为季度固定效应。问题在于
$\text{tone}_t$ 在同一季度内对所有行业取同一个值，它是 $t$ 的函数；而
$\lambda_t$ 已经为每个季度提供了一个独立的截距。二者完全共线，$\beta$ 无法识别，
估计软件会直接把 $\text{tone}_t$ 标记为共线并剔除。

这不是数据量不足或样本期太短造成的，而是模型设定与变量层级的必然结果：任何只在时间
维度上变化的解释变量，在加入时间固定效应后都不可识别。行业面板中确实存在大量
$i\times t$ 层面的变异，但情绪没有这一层面的变异。

由此产生两条可行的路径。第一条是放弃时间固定效应，直接在时间序列上识别，代价是无法
控制同期的共同冲击；第二条是保留固定效应，转而估计情绪与行业特征的交互项
$\text{tone}_t\times Z_i$，此时交互项在 $i\times t$ 层面有变异，斜率可以识别。本文采用
第一条路径，理由在下一小节说明；第二条路径作为后续扩展在 4.5 节讨论。

\subsection{以事件窗口作为识别来源}

本文把识别建立在报告发布日附近的市场反应上，而不是季度的平均状态。这样做的考虑有三点。

其一，报告的影响在时间上是集中的。政策信息在发布当天进入价格，若用整个季度的收益去
回归，等于把事件效应与季度内的其他信息混在一起，稀释了信号。事件窗把观测限制在信息
到达的附近，噪音相对可控。

其二，事件窗内的比较是“同一时点、同一批行业”，共同冲击被时间维度自动吸收。所有行业
在同一交易日面对同一份报告，行业间的差异不需要靠固定效应剔除，样本本身已经是同日
对照。

其三，事件设计使“报告期”与“发布期”的错位得到处理。第 3 章已经说明，报告总是在报告期
结束后的下一季度发布；以发布日为事件起点，可以避免用事后信息解释事前行情。

需要明确的是，这一设计识别的是发布日前后的条件关联，而非严格的因果效应。报告内容与
当期经济状况相关，措辞的变化可能同时反映基本面，因此系数的解释应限定为“给定发布时点，
情绪较高与随后的市场反应之间的关系”。本文在 4.5 节讨论这一限制及其缓解办法。

\subsection{基准模型}

对每一期报告 $e$，估计
\begin{equation}
Y_e=\alpha+\beta\,\text{tone}_{P,e}+\gamma'X_e+\varepsilon_e,
\label{eq:baseline}
\end{equation}
其中 $Y_e$ 为事件窗内的市场收益或波动指标，取一日收益 $R^{(1)}$、五日收益 $R^{(5)}$、
三日累计收益 CAR、事件日绝对收益 $|R^{(0)}|$ 与日内振幅 Range 五种；
$\text{tone}_{P,e}$ 为标准化后的政策指引段情绪；$X_e$ 为控制变量。

控制变量的选择针对一个具体的混淆。情绪较高的报告往往出现在市场相对平静的时期，
若不加控制，情绪与低波动之间的相关可能只是“平静期写得更乐观”的反映。为此在波动
方程中加入事前 60 个交易日的市场收益标准差 $\text{PreVol}_e$，用以吸收波动聚集。
在双段回归中同时纳入宏观回顾段情绪 $\text{tone}_{R,e}$，用于区分“政策指引”与
“整篇语气”两种解释。

标准误采用 Newey--West 估计，滞后阶数取 4。事件间隔约一个季度，事件窗之间不重叠，
因此主要考虑的是残差的自相关与异方差，而非重叠窗口带来的相关性。

与既有研究常用的 EGARCH 设定相比，本文选择事件窗回归出于两点考虑。一是事件窗直接
对应信息到达的时点，条件方差模型则把波动视为状态的连续演化，难以把效应归到某一期
报告上。二是本文需要检验断点，事件窗回归可以通过交互项直接完成，而 GARCH 族模型的
参数时变性需要另作设定。为弥补这一差异，本文同时使用绝对收益与日内振幅两个波动代理，
并在稳健性中报告五日窗口的结果。

\subsection{制度断点的设定}

第 5 章关心的第二个问题是效应是否稳定。设定的出发点是样本期内货币政策框架的调整：
2013 年的流动性紧张、2014 年利率市场化的推进、2015 年汇率形成机制的调整，以及报告
自身章节结构在 2015 年前后的改写。这些变化都可能改变市场解读文本的方式。

据此在式~\eqref{eq:baseline} 中加入断点虚拟变量及其交互项：
\begin{equation}
Y_e=\alpha+\beta_1\,\text{tone}_{P,e}
+\beta_2\,\text{tone}_{P,e}\times\text{post}_{2014}
+\beta_3\,\text{post}_{2014}+\varepsilon_e,
\label{eq:break}
\end{equation}
其中 $\text{post}_{2014}$ 在 2014 年及以后取 1。$\beta_1$ 为断点前的斜率，
$\beta_1+\beta_2$ 为断点后的斜率，$\beta_2$ 即效应的变化量。

断点年份由制度节点确定，而非在数据中搜索得到。为说明这一点，本文把断点依次替换为
2008、2010、2012、2014、2016、2018 年重新估计，报告交互项系数与 $t$ 值的变化。
若结论只在某一年份成立，则说明断点是人为挑选的；若多个相邻年份方向一致，则说明变化
是渐进的。第 5 章表 5.4 报告了这一结果。

\subsection{稳健性设计与识别假设}

\paragraph{四类稳健性检验}
第一，更换市场口径。把等权组合由 16 个行业扩展为 28 个行业，检验结论是否依赖行业集合。
第二，更换情绪词典。以自建货币政策领域词典重新度量情绪，观察系数方向与显著性如何变化。
第三，影响点检验。逐一剔除 2015 年（股市异常波动）、2020 年（疫情冲击）与 2003Q1
（分段异常），并对收益作 1\%/99\% 缩尾，检验结论是否由少数观测支撑。
第四，置换检验。将情绪与事件随机配对 500 次，统计真实 $t$ 值在随机分布中的位置，
这一检验不依赖残差分布的正态假设。

\paragraph{统计功效}
在解读“不显著”之前，需要知道这条设计能识别多大的效应。以一日收益方程为例，系数标准误
约为 25 个基点；在 5\% 显著性水平、80\% 功效下，可识别的最小效应约为 70 个基点，
而日频市场收益本身的标准差接近 190 个基点。这意味着小于 70 个基点的效应在本样本中
无法被可靠识别，不显著的系数不能直接解读为“效应为零”。本文在报告收益方程结果时
一并说明这一限制。

\paragraph{尚不能排除的问题}
其一，同期公告的干扰。货币政策操作、金融统计数据、物价数据等常在相近时点发布，
事件窗内可能同时包含其他信息。本文未能逐日控制这些公告，只能通过窗口长度与安慰剂
检验间接缓解。

其二，发布时点不明。元数据中不含发布的具体时刻，无法判断报告是在盘中还是收盘后
发布。本文以发布当日为 $t_0$；由于无法判断确切时刻，第 5 章同时报告三日累计收益与
五日窗口的结果，用以检验结论对窗口起点的敏感性。

其三，情绪度量的误差。词典法只能识别词面极性，无法处理反讽、条件句与语境依赖；
测量误差会使系数向零收缩，因而本文报告的效应量级更可能是下界。第 5 章以替代词典
检验说明这一偏误的方向与幅度。

其四，反向因果。报告措辞可能与当期经济状况共同决定，本文的断点与置换检验只能缓解
而不能排除这一可能。因此全文的表述避免使用“导致”，而采用“与……相关”“显著降低”
这类描述性措辞。

\paragraph{关于行业层面}
本文的市场组合由行业等权构造，28 行业口径用作稳健性，但并未单独估计“情绪对哪些行业
影响更大”。原因在于 4.1 节所述的识别约束：在季度面板中，情绪与时间固定效应共线，
即使加入行业特征交互项，也需要行业特征本身在样本期内有足够变异，而现有数据中可用的
行业特征（估值、杠杆、利率敏感度）覆盖期偏短，难以支撑稳定估计。这一扩展留待后续
研究，本文对此不作结论。
